% 2.1.7
% Ficheiro incluído por main.tex

\subsection[Login com contas do AD: PAM]{Login com contas do AD: PAM}

\esquema{2-1-7-pam}{Login de um utilizador do AD com PAM e SSSD}

\passo{Passo 1 - Ativar a criação automática da home no primeiro login}

\begin{lstlisting}[style=bashstyle]
sudo pam-auth-update --enable mkhomedir
grep mkhomedir /etc/pam.d/common-session
\end{lstlisting}

\emph{Verificação esperada:} o \texttt{grep} mostra \texttt{pam\_mkhomedir.so}. Não há MFA configurado: o login pede só utilizador e palavra-passe.

\passo{Passo 2 - Login com uma conta do AD}

\begin{lstlisting}[style=bashstyle]
su - aluno.teste@isec.local
whoami
id
pwd
\end{lstlisting}

\emph{Verificação esperada:} o \texttt{su -} autentica o aluno no AD através do \textbf{PAM + SSSD}; \texttt{whoami} mostra \texttt{aluno.teste@isec.local}; a home é criada na primeira vez (\texttt{pwd} → \texttt{/\allowbreak{}home/\allowbreak{}aluno.teste@isec.local}).

\begin{importantebox}[title={Importante}]
\textbf{Se o \texttt{su -} der \texttt{su: System error}}, o SSSD está a aplicar as GPO do AD ao login e bloqueia-o. Volta ao terminal do utilizador local e corrige o \texttt{sssd.conf}:

\begin{lstlisting}[style=bashstyle]
# [Linux Client]  (terminal do utilizador local)
sudo nano /etc/sssd/sssd.conf
\end{lstlisting}

Na secção \texttt{[domain/\allowbreak{}isec.local]}, acrescenta a linha \texttt{ad\_gpo\_access\_control = permissive} (e apaga \texttt{config\_file\_version = 2}, se existir). Guarda (\texttt{Ctrl+O}, \texttt{Enter}, \texttt{Ctrl+X}) e reinicia o SSSD:

\begin{lstlisting}[style=bashstyle]
sudo chmod 600 /etc/sssd/sssd.conf
sudo sssctl config-check
sudo systemctl restart sssd
\end{lstlisting}

Depois repete este passo. O modo \texttt{permissive} regista as GPO mas não as impõe (solução de laboratório).
\end{importantebox}

\passo{Passo 3 - Obter o ticket Kerberos na sessão}

\begin{lstlisting}[style=bashstyle]
kinit aluno.teste@ISEC.LOCAL
klist
\end{lstlisting}

\emph{Verificação esperada:} \texttt{klist} mostra \texttt{krbtgt/\allowbreak{}ISEC.LOCAL@ISEC.LOCAL}.

\emph{Porquê:} o ticket pertence a quem o pede - um \texttt{kinit} feito como utilizador local não serve ao aluno. Por isso faz-se \textbf{dentro} da sessão do aluno. \textbf{Não saias da sessão:} continua em 2.1.8.

\problemas{\hyperref[ec-b1-pam]{Erros Comuns - 2.1 PAM e pam\_mount}}
